% Generated draft: requires mathematical and source review.
\documentclass[11pt]{article}
\usepackage[margin=25mm]{geometry}
\usepackage{fontspec}
\setmainfont{Latin Modern Roman}
\usepackage{amsmath,amssymb,mathtools,unicode-math}
\setmathfont{Latin Modern Math}
\usepackage{enumitem,needspace,xurl,hyperref}
\setlength{\parindent}{0pt}
\setlength{\parskip}{5pt}
\setlength{\emergencystretch}{3em}
\newcommand{\N}{\mathbb{N}}
\newcommand{\Z}{\mathbb{Z}}
\newcommand{\Q}{\mathbb{Q}}
\newcommand{\R}{\mathbb{R}}
\newcommand{\C}{\mathbb{C}}
\newcommand{\F}{\mathbb{F}}
\newcommand{\A}{\mathbb{A}}
\newcommand{\PP}{\mathbb{P}}
\newcommand{\E}{\mathbb{E}}
\newcommand{\eps}{\varepsilon}
\newcommand{\dd}{\,\mathrm{d}}
\begin{document}
\section*{UnsolvedMath ID: 30004306}
\textbf{Source problem code:} catalogue entry
\subsection{Definitions and mathematical statement}
The supplied question is: \emph{Is there an algebraic invariant that determines the order of the Yang--Baxter map \(r\)?} The record does not specify the domain of \(r\), the class of Yang--Baxter maps under consideration, or what qualifies as an algebraic invariant.

For purposes of interpreting the word \emph{order}, if \(r\colon S\to S\) is a bijection of a set \(S\), define
\[
\operatorname{ord}(r)=
\begin{cases}
\min\{n\in\Z_{\geq 1}:r^n=\operatorname{id}_S\},&\text{if this set is nonempty},\\
\infty,&\text{otherwise}.
\end{cases}
\]
Here \(r^n\) denotes composition of \(r\) with itself \(n\) times. Bijectivity is a condition for this permutation-order interpretation; it is not explicitly supplied in the record.

One possible setting is a set-theoretic solution in the braid convention: a set \(X\) and a bijection \(r\colon X\times X\to X\times X\) satisfying
\[
(r\times\operatorname{id}_X)\circ(\operatorname{id}_X\times r)\circ(r\times\operatorname{id}_X)
=
(\operatorname{id}_X\times r)\circ(r\times\operatorname{id}_X)\circ(\operatorname{id}_X\times r)
\quad\text{on }X^3.
\]
This is an illustrative interpretation, not a hypothesis established by the supplied source record. The intended convention for the Yang--Baxter equation must be confirmed, together with any finiteness, nondegeneracy, or involutivity assumptions.

Once an intended class \(\mathcal C\) and an equivalence relation on its objects have been specified, an invariant would be an assignment \(I\colon\mathcal C\to A\) constant on equivalent objects. To say that it determines the order means
\[
\forall u,v\in\mathcal C,\qquad
I(u)=I(v)\ \Longrightarrow\ \operatorname{ord}(r_u)=\operatorname{ord}(r_v),
\]
where \(r_u\) is the Yang--Baxter map belonging to \(u\). This gives a conditional meaning to ``determines''; neither \(\mathcal C\), the equivalence relation, nor the permitted algebraic construction of \(I\) is supplied. In particular, permitting the order itself as the invariant would make the question tautological. A precise substantive existence problem therefore requires additional restrictions on the intended invariant. The supplied question is retained with these qualifications unresolved.
\subsection{Short English statement}
The question asks whether algebraic data associated with a Yang--Baxter map determine its order under iteration. The intended maps, hypotheses, and permitted invariants are unspecified, so the displayed formulation is conditional and requires mathematical review. The supplied record labels the problem open; this status has not been independently verified.
\subsection{Sources}
\begin{itemize}
\item Local UnsolvedMath registry, permanent ID \(30004306\), title: \emph{Algebraic Invariants of Yang--Baxter Map Order}; collection: Oberwolfach Reports Open Problems; exact source problem code: \texttt{catalogue entry}. No individual proposer is supplied. The proposed year is recorded as 2019.
\item Supplied source citation: \emph{Mini-Workshop: Algebraic Tools for Solving the Yang--Baxter Equation (2020)}, Oberwolfach Reports, pp.~3207--3242, DOI 10.4171/owr/2019/51. \url{https://doi.org/10.4171/owr/2019/51}. The record describes Oberwolfach Reports as documenting mathematical workshops and their open-problem discussions. The supplied title year 2020 and the recorded proposed year 2019 are retained without resolving their relationship.
\item Supplied literature reference, labelled ``Gateva--Ivanova, derived symmetric and permutation groups'': \url{https://arxiv.org/abs/1507.02602}.
\item Supplied literature reference, labelled ``original OWR report'': \url{https://doi.org/10.4171/owr/2019/51}.
\item Literature assessment dated exactly \texttt{2026-08-22}: the supplied curation labels the problem open and reports that derived permutation groups, retractions, and related invariants provide partial information about Yang--Baxter solutions, while no general algebraic invariant determining the order of \(r\) has been established. This is supplied dated curation derived from the OWR record, not an independently verified literature assessment. Its time-sensitive status and cited claims require independent verification.
\item The registry labels the statement \texttt{corrected\_verified} and reports a manual semantic, typographic, and LaTeX review asserting that the clean statement is self-contained and preserves the target. These are supplied editorial assertions; they do not resolve the formulation issues identified here.
\end{itemize}
\end{document}
